\subsubsection{自适应辛普森 (Adaptive Simpson's Rule)}
自适应 Simpson，传入函数与容差；复用端点／中点值以减少函数求值。每次调用使用局部状态，多测须重新捕获本组函数参数。区间内须连续且足够光滑，不可含极点或未经处理的广义积分；给出绝对误差估计，不是对任意函数的严格误差保证。
\begin{lstlisting}
template<class F>
double integrate(F func,double left,double right,double eps=1e-8,int max_depth=30){
    assert(eps>0&&max_depth>0);
    if(left==right)return 0;
    if(left>right)return -integrate(func,right,left,eps,max_depth);
    auto eval=[&](double x){double y=func(x);if(!isfinite(y))throw domain_error("non-finite integrand");return y;};
    auto simpson=[](double l,double r,double fl,double fm,double fr){return (r-l)*(fl+4*fm+fr)/6;};
    function<double(double,double,double,double,double,double,double,int)> rec;
    rec=[&](double l,double r,double fl,double fm,double fr,double old,double tol,int depth)->double{
        double m=l+(r-l)/2;
        if(m==l||m==r)return old;
        double lm=eval(l+(m-l)/2),rm=eval(m+(r-m)/2);
        double a=simpson(l,m,fl,lm,fm),b=simpson(m,r,fm,rm,fr),delta=a+b-old;
        if(fabs(delta)<=15*tol)return a+b+delta/15;
        if(!depth)throw runtime_error("Simpson did not converge");
        return rec(l,m,fl,lm,fm,a,tol/2,depth-1)+rec(m,r,fm,rm,fr,b,tol/2,depth-1);
    };
    double m=left+(right-left)/2,fl=eval(left),fm=eval(m),fr=eval(right);
    return rec(left,right,fl,fm,fr,simpson(left,right,fl,fm,fr),eps,max_depth);
}
\end{lstlisting}
